% !PRINTED 497
\weberpage{497}
%<<CONT
erkennt, und daraus folgt, dass eine cyklische Permutation zur ersten oder zur zweiten Art gehört, je nachdem die Anzahl der Ziffern eine ungerade oder eine gerade ist.

Wie man jede Permutation aus Transpositionen zusammensetzen kann, ebenso kann man jede Permutation der ersten Art aus dreigliedrigen Cyklen zusammensetzen. Es genügt, wenn dies für jedes Paar von Transpositionen bewiesen ist, da jede Permutation der ersten Art sich aus solchen Paaren componiren lässt. Es ist aber
\[
(0, 1)(0, 2) = (0, 1, 2)
\]
\[
(0, 3)(1, 2) = (0, 1, 2)(0, 1, 3),
\]
woraus die Richtigkeit der Behauptung zu ersehen ist, da die beiden Transpositionen eines Paares entweder eine oder keine Ziffer gemein haben. Also:

\gesperrt{6. Alle Permutationen der ersten Art lassen sich (auf unendlich viele Weisen) in cyklische Permutationen von drei Elementen zerlegen.}

Aus 6. folgt weiter

\gesperrt{7. Eine Permutationsgruppe, die alle dreigliedrigen cyklischen Permutationen mit zwei festen Ziffern enthält, muss die ganze alternirende Gruppe enthalten.}

Denn aus den Cyklen $(0, 1, 2), (0, 1, 3), (0, 1, 4) \dots (0, 1, m - 1)$ lassen sich alle dreigliedrigen Cyklen und also nach 6. alle Permutationen der ersten Art componiren, wie man aus den Zusammensetzungen
\[
(1, 0, 2) = (0, 1, 2)(0, 1, 2),
\]
\[
(0, 2, 3) = (0, 1, 2)(1, 0, 3)(1, 0, 2),
\]
\[
(1, 3, 4) = (1, 0, 3)(0, 1, 4)(0, 1, 3)
\]
\[
(2, 3, 4) = (1, 0, 2)(1, 3, 4)(0, 1, 2)
\]
ersieht.

\gesperrt{8. Ist $m$ grösser als 4, so ist eine Permutationsgruppe, die alle aus je zwei Transpositionen ohne gemeinsame Ziffer zusammengesetzten Permutationen $(0, 1)(2, 3)$ enthält, durch die alternirende Gruppe theilbar.}

Denn es ist
\[
(0, 1, 2) = (0, 1)(3, 4)(3, 4)(0, 2)
\]
%CONT>>
